% ============================================================
% Section 4C: Rigorous Structural Enhancement Framework
% Revised based on Codex (OpenAI GPT-5.2) strict review
% Key changes:
%   1. Acknowledge we study "marginal gain with side information"
%   2. Use correct sigma-algebra definitions
%   3. Distinguish Bayes-optimal from algorithm-achievable
%   4. Honest about limitations
% ============================================================

\subsection{Structural Enhancement with Side Information}
\label{subsec:structural_enhancement_revised}

We now develop a rigorous framework for understanding when graph structure provides additional benefit beyond node features. Following the critique that standard Kesten-Stigum results apply to ``pure structure'' detection, we carefully distinguish our setting.

\subsubsection{Precise Problem Formulation}

\begin{definition}[Contextual SBM with Gaussian Features]
\label{def:csbm_precise}
Fix $n \in \mathbb{N}$ and parameters $(\rho, d, \text{SNR}_X)$ where:
\begin{itemize}
    \item $\rho \in [-1, 1]$: Label correlation across edges
    \item $d > 0$: Expected degree (Poisson$(d)$ model)
    \item $\text{SNR}_X = \|\mu\|^2 / \sigma^2$: Feature signal-to-noise ratio
\end{itemize}

Generate data as follows:
\begin{enumerate}
    \item Labels: $Y_v \stackrel{\text{iid}}{\sim} \mathrm{Rademacher}(\frac{1}{2})$ for $v \in [n]$
    \item Features: $X_v = \mu Y_v + \varepsilon_v$ where $\varepsilon_v \stackrel{\text{iid}}{\sim} \mathcal{N}(0, \sigma^2 I_d)$
    \item Edges: For each unordered pair $\{u, v\}$, independently include edge with probability
    \begin{equation}
        p_{uv} = \frac{d}{n} \cdot \left(1 + \rho \cdot Y_u Y_v\right) / 2
    \end{equation}
\end{enumerate}
The homophily parameter is $h = (1 + \rho)/2$, giving $\rho = 2h - 1$.
\end{definition}

\begin{definition}[Information Sets]
\label{def:info_sets}
For node $v \in [n]$, define $\sigma$-algebras:
\begin{itemize}
    \item $\mathcal{F}_0(v) := \sigma(X_v)$: Own features only
    \item $\mathcal{F}_k(v) := \sigma(X_v, G_k(v), \{X_u : u \in B_k(v)\})$: Own features, $k$-hop induced subgraph, and all features within $k$-hop neighborhood
\end{itemize}
where $B_k(v) = \{u : d_G(u, v) \leq k\}$ is the $k$-hop ball and $G_k(v)$ is the induced subgraph on $B_k(v)$.
\end{definition}

\begin{definition}[Classification Accuracies]
\label{def:accuracies_precise}
\begin{itemize}
    \item \textbf{MLP accuracy}: $\mathrm{Acc}_{\mathrm{MLP}}^* := \mathbb{E}[\mathbf{1}\{\hat{Y}_v^{\mathrm{MLP}} = Y_v\}]$ where $\hat{Y}_v^{\mathrm{MLP}} = \mathrm{sign}(\langle \mu, X_v \rangle)$ is the Bayes-optimal classifier using $\mathcal{F}_0(v)$
    \item \textbf{$k$-hop Bayes accuracy}: $\mathrm{Acc}_{k}^* := \mathbb{E}[\mathbf{1}\{\hat{Y}_v^{(k)} = Y_v\}]$ where $\hat{Y}_v^{(k)} = \mathbb{E}[Y_v \mid \mathcal{F}_k(v)]$ is the Bayes-optimal classifier using $\mathcal{F}_k(v)$
    \item \textbf{Marginal gain}: $\Delta_k := \mathrm{Acc}_{k}^* - \mathrm{Acc}_{\mathrm{MLP}}^*$
\end{itemize}
\end{definition}

\subsubsection{The Marginal Gain Problem}

Our central question is: \textit{When is $\Delta_k > 0$, and how large can it be?}

This differs from the classical KS problem (``can we detect communities from $G$ alone?'') because:
\begin{enumerate}
    \item We already have side information $X_v$ that partially reveals $Y_v$
    \item We ask about \textit{marginal} improvement, not absolute detectability
    \item The relevant mutual information is $I(Y_v; \mathcal{F}_k(v) \mid X_v)$, not $I(Y_v; G)$
\end{enumerate}

\subsubsection{Fundamental Bounds}

\begin{theorem}[Upper Bound on Marginal Gain]
\label{thm:marginal_gain_upper}
For any $k \geq 1$:
\begin{equation}
    \Delta_k \leq \mathrm{Acc}_{k}^* - \mathrm{Acc}_{\mathrm{MLP}}^* \leq 1 - \mathrm{Acc}_{\mathrm{MLP}}^* =: \mathcal{B}
\end{equation}
The Information Budget $\mathcal{B}$ is the maximum possible gain.
\end{theorem}

\begin{proof}
Since $\mathrm{Acc}_{k}^* \leq 1$ and $\mathrm{Acc}_{\mathrm{MLP}}^*$ is fixed, the bound follows immediately.
\end{proof}

\begin{theorem}[Conditional Mutual Information Bound]
\label{thm:cmi_bound}
For the marginal gain:
\begin{equation}
    \Delta_k \leq \sqrt{\frac{\ln 2}{2} \cdot I(Y_v; \mathcal{F}_k(v) \mid X_v)}
\end{equation}
where the right-hand side is the Pinsker bound applied to the conditional distribution.
\end{theorem}

\begin{proof}
By the relation between accuracy and total variation, and Pinsker's inequality applied to $P(Y_v \mid \mathcal{F}_k(v))$ vs $P(Y_v \mid X_v)$.
\end{proof}

\subsubsection{Characterization of Marginal Gain (Empirical)}

Due to the complexity of analyzing $I(Y_v; \mathcal{F}_k(v) \mid X_v)$ with both structure and neighbor features, we provide an \textbf{empirical characterization} rather than a full proof.

\begin{observation}[Structural SNR and Marginal Gain]
\label{obs:snr_marginal}
Define the \textbf{structural SNR} as $\mathrm{SNR}_G := (d-1)\rho^2$.

Experimentally, we observe:
\begin{enumerate}
    \item When $\mathrm{SNR}_G \leq 1$: Marginal gain $\Delta_k \approx 0$ for all $k$
    \item When $\mathrm{SNR}_G > 1$: Marginal gain $\Delta_k > 0$ for appropriate $k$
\end{enumerate}
This suggests the KS threshold may govern marginal gains, but a formal proof requires analyzing ``broadcasting with side information.''
\end{observation}

\begin{remark}[Open Problem: KS Threshold with Side Information]
\label{rem:open_problem}
The precise characterization of $\Delta_k$ in terms of $(\mathrm{SNR}_X, \mathrm{SNR}_G, k)$ remains open. Key technical challenges:
\begin{enumerate}
    \item \textbf{Coupled state evolution}: BP with feature-initialized beliefs has different dynamics than uninformed BP
    \item \textbf{Information-theoretic vs algorithmic}: Even if $I(Y_v; \mathcal{F}_k(v) \mid X_v) > 0$, it may not be extractable by efficient algorithms
    \item \textbf{Finite-$n$ effects}: Local weak convergence arguments require careful control of cycle structure
\end{enumerate}
We leave the full theoretical treatment to future work.
\end{remark}

\subsubsection{Practical Phase Diagram}

Despite theoretical gaps, our experiments support a clear \textbf{practical phase diagram}:

\begin{table}[t]
\centering
\caption{Empirical Phase Diagram for GNN vs MLP}
\label{tab:empirical_phase}
\begin{tabular}{@{}lccl@{}}
\toprule
\textbf{Regime} & \textbf{Condition} & \textbf{Marginal Gain} & \textbf{Recommendation} \\
\midrule
Feature-sufficient & $\mathcal{B} < 0.05$ & $\approx 0$ & MLP \\
Structure-beneficial (high $h$) & $h > 0.7$ and $\mathcal{B} > 0.05$ & $> 0$ & GNN \\
Structure-harmful (mid $h$) & $0.3 < h < 0.7$ & $< 0$ (ADR) & MLP or Concat-GNN \\
Heterophily-recoverable & $h < 0.3$ and $R > 1.5$ & $> 0$ & H2GCN \\
Heterophily-irrecoverable & $h < 0.3$ and $R \leq 1$ & $\approx 0$ & MLP \\
\bottomrule
\end{tabular}
\end{table}

\begin{remark}[Empirical vs Theoretical]
The thresholds $(0.3, 0.7)$ and $R = 1.5$ are \textbf{empirically calibrated} from our experiments, not theoretically derived. They approximately correspond to $\mathrm{SNR}_G = 1$ for graphs with average degree $d \approx 7$--$15$.

We emphasize: \textbf{this is an empirical framework with theoretical motivation}, not a fully proven theory. The practical value lies in predictive accuracy (88.9\% on CSBM, 77.8\% external validation), not mathematical certainty.
\end{remark}

\subsubsection{Honest Statement of Theoretical Contribution}

Our theoretical contributions are:

\begin{enumerate}
    \item \textbf{Information Budget Bound} (Theorem~\ref{thm:budget_pinsker}): Rigorous upper bound on GNN improvement using Pinsker's inequality

    \item \textbf{Shrinkage Factor} (Theorem~\ref{thm:discriminability_shrinkage}): Exact formula $|2h-1|$ for class separation under mean aggregation

    \item \textbf{Concatenation Preservation} (Theorem~\ref{thm:concat_preserves}): DPI-based guarantee that concatenation cannot hurt

    \item \textbf{KS-Motivated Phase Diagram} (Observation~\ref{obs:snr_marginal}): Empirical evidence that $\mathrm{SNR}_G = 1$ predicts regime transitions
\end{enumerate}

The first three are \textbf{mathematically rigorous}. The fourth is \textbf{empirically supported but theoretically incomplete}. We acknowledge this limitation and propose it as direction for future theoretical work.

\begin{remark}[Comparison with Pure KS Results]
Our setting differs from classical KS in a crucial way: we ask ``does structure help \textit{given features}?'' rather than ``can structure alone detect communities?'' This ``marginal gain with side information'' problem is less studied and harder to analyze. Our empirical finding that the same threshold $(d-1)\rho^2 = 1$ appears relevant is intriguing but unproven.
\end{remark}


